\chapter{Lista tuturor constrângerilor ce trebuie îndeplinite de model}
\label{sec:constrangeri}

\section*{a. Unicitate}

\textbf{i. Unicitate locală}\\
\begin{itemize}
    \item Câmpul \textit{cnp} din fragmentul \textit{CLIENT\_FIZICA} pe \textbf{SV1} trebuie să fie unic pentru a identifica persoanele fizice.
    \item Câmpul \textit{cui} din fragmentul \textit{CLIENT\_JURIDICA} pe \textbf{SV2} trebuie să fie unic pentru a identifica entitățile juridice.
\end{itemize}

\textbf{ii. Unicitate globală pe fragmente orizontale}\\
\begin{itemize}
    \item Câmpul \textit{email} din tabela \textit{CLIENT} trebuie să fie unic la nivelul întregului sistem. Deoarece clienții sunt distribuiți între \textbf{SV1} (B2C) și \textbf{SV2} (B2B), unicitatea va fi asigurată prin interogarea ambelor noduri prin \textit{Database Links} înaintea inserării.
\end{itemize}

\textbf{iii. Unicitate globală pe fragmente verticale}\\
\begin{itemize}
    \item Câmpul \textit{agent\_id} asigură legătura unică între datele publice din \textit{AGENT\_PROFIL} (\textbf{SV1}/\textbf{SV2}) și datele sensibile din \textit{AGENT\_SEC} (\textbf{SV3}). Această unicitate este critică pentru securitatea procesului de autentificare.
\end{itemize}


\section*{b. Cheie Primară (la nivel local/global)}

\begin{itemize}
    \item La nivel \textbf{local}, tabelele precum \textit{ADRESA} sau \textit{TICKET\_HISTORY} folosesc coloane de tip \textit{IDENTITY} pentru a genera chei primare unice pe nodul respectiv.
    \item La nivel \textbf{global}, pentru a evita coliziunile între \textit{client\_id} generate pe \textbf{SV1} și \textbf{SV2}, se vor defini intervale de valori (Sequences) disjuncte (ex: \textbf{SV1} folosește $1 \dots 10^6$, \textbf{SV2} folosește $10^6+1 \dots 2 \cdot 10^6$).
\end{itemize}


\section*{c. Cheie Externă (la nivel local și global)}

\textbf{i. Chei externe locale}\\
\begin{itemize}
    \item \textit{TICKET\_FIZIC.client\_id} face referință la \textit{CLIENT\_FIZICA.client\_id} pe \textbf{SV1}.
    \item \textit{ADRESA.client\_id} face referință la clienții locali de pe nodurile respective.
\end{itemize}

\textbf{ii. Chei externe globale (Cross-DB)}\\
\begin{itemize}
    \item \textit{TICKET.status\_id} face referință la tabela \textit{STATUS} aflată pe \textbf{SV4}. Deoarece \textit{STATUS} este replicată pe \textbf{SV1} și \textbf{SV2} ca \textit{Materialized View}, integritatea este verificată local, dar sursa de adevăr rămâne nodul central.
\end{itemize}


\section*{d. Validare (la nivel local și global)}

\textbf{i. Validări locale (Check Constraints)}\\
\begin{itemize}
    \item \textit{client\_type} trebuie să fie 'F' sau 'J'.
    \item \textit{este\_final} din tabela \textit{STATUS} trebuie să fie 'Y' sau 'N'.
    \item \textit{nivel} de prioritate trebuie să fie între 1 și 5.
\end{itemize}

\textbf{ii. Validări globale}\\
\begin{itemize}
    \item Un tichet nu poate fi atribuit unui agent dacă statusul \textit{is\_active} al agentului (stocat în \textit{AGENT\_PROFIL} pe \textbf{SV1}/\textbf{SV2}) este 'N'.
\end{itemize}